\label{sec:alcance}
Se estudiará un mecanismo sobre una familia PIKAN seleccionada mediante piloto. B-splines y centros compactos permanecerán como candidatos hasta verificar representación, derivadas y costo. Se utilizarán PDEs de baja dimensión y un sistema piezoeléctrico lineal estacionario. El objetivo es caracterizar un mecanismo, no demostrar superioridad universal de una arquitectura.

El núcleo confirmatorio incluye Poisson localizado y Burgers; el banco distribuido y las señales permutadas controlan explicaciones alternativas. Allen--Cahn por ventanas evalúa seguimiento, Navier--Stokes verifica acoplamiento y geometría, y piezoelectricidad evalúa transmisión y reparto entre campos. El caso dimensional PVDF requiere ficha de material verificada; mientras tanto se dispone de un caso sintético completamente especificado.

Quedan fuera turbulencia, grandes deformaciones, validación experimental de dispositivos y afirmaciones de novedad absoluta. Las referencias FEM se someterán a convergencia. Las figuras analíticas y la demostración FEM del protocolo no constituyen evidencia de eficacia del controlador. El alcance y los umbrales propuestos requieren revisión científica antes de activar la campaña.
